In industries such as airlines, steel, and truck rentals, 
firms have coordinated to raise prices or reduce capacities 
using public announcements such as earnings calls. As this 
type of coordinating practice is relatively unexplored, this study seeks 
to learn more about the collusive content 
in firms' public communications. To do so, we design a 
large language model (LLM) to look for such content and 
then apply it to almost 500,000 company transcripts. The 
LLM is very selective in that only .12\% of the transcripts 
satisfy the criterion for collusive content. The LLM is 
reasonably effective at rediscovering previously documented 
episodes and identifies many new episodes. A human audit 
is conducted of the 301 transcripts with the highest LLM 
scores. The audit distills and organizes the content to 
identify systematic ways in which it is facilitating 
coordinated conduct. Messages fall into two broad bins. 
First, a firm states that the industry needs to raise 
prices, reduce supply, or reduce capacity. Second, a 
firm expresses how it has been less aggressive and 
calls on competitors to join them, either explicitly 
or implicitly by saying it is "acting as a leader", 
commending rival firms who have acted in a similar 
manner, or criticizing rival firms that have not. A 
second research objective is to develop a screening 
tool for competition agencies and plaintiff law firms. 
Towards that end, the human audit evaluates the 
performance of the LLM in substituting for human experts. 
The LLM is capable of identifying a wide variety of 
collusive content though there is a reasonably high 
false positive rate. Nevertheless, we explain how it 
can be an essential component of a screening protocol 
and show proof of concept with a case study of the 
Brazilian domestic airline industry.
